\documentclass[10pt,a4paper]{amsart}
\usepackage{amsmath,amssymb,amsthm,graphicx}
\usepackage{verbatim}
\newenvironment{reference}{\comment}{\endcomment}
%\newenvironment{reference}{}{}
\newenvironment{slogan}{\comment}{\endcomment}
\newenvironment{history}{\comment}{\endcomment}

% For commutative diagrams we use Xy-pic
\usepackage[all]{xy}

% We use 2cell for 2-commutative diagrams.
\xyoption{2cell}
% Theorem environments.
%
\theoremstyle{plain}
\newtheorem{theorem}[subsection]{Theorem}
\newtheorem{proposition}[subsection]{Proposition}
\newtheorem{lemma}[subsection]{Lemma}

\theoremstyle{definition}
\newtheorem{definition}[subsection]{Definition}
\newtheorem{example}[subsection]{Example}
\newtheorem{exercise}[subsection]{Exercise}
\newtheorem{situation}[subsection]{Situation}

\theoremstyle{remark}
\newtheorem{remark}[subsection]{Remark}
\newtheorem{remarks}[subsection]{Remarks}

\numberwithin{equation}{subsection}

% Macros
%
\def\lim{\mathop{\mathrm{lim}}\nolimits}
\def\colim{\mathop{\mathrm{colim}}\nolimits}
\def\Spec{\mathop{\mathrm{Spec}}}
\def\Hom{\mathop{\mathrm{Hom}}\nolimits}
\def\Ext{\mathop{\mathrm{Ext}}\nolimits}
\def\SheafHom{\mathop{\mathcal{H}\!\mathit{om}}\nolimits}
\def\SheafExt{\mathop{\mathcal{E}\!\mathit{xt}}\nolimits}
\def\Sch{\mathit{Sch}}
\def\Mor{\mathop{\mathrm{Mor}}\nolimits}
\def\Ob{\mathop{\mathrm{Ob}}\nolimits}
\def\Sh{\mathop{\mathit{Sh}}\nolimits}
\def\NL{\mathop{N\!L}\nolimits}
\def\CH{\mathop{\mathrm{CH}}\nolimits}
\def\proetale{{pro\text{-}\acute{e}tale}}
\def\etale{{\acute{e}tale}}
\def\QCoh{\mathit{QCoh}}
\def\Ker{\mathop{\mathrm{Ker}}}
\def\Im{\mathop{\mathrm{Im}}}
\def\Coker{\mathop{\mathrm{Coker}}}
\def\Coim{\mathop{\mathrm{Coim}}}

% Boxtimes
%
\DeclareMathSymbol{\boxtimes}{\mathbin}{AMSa}{"02}

%
% Macros for moduli stacks/spaces
%
\def\QCohstack{\mathcal{QC}\!\mathit{oh}}
\def\Cohstack{\mathcal{C}\!\mathit{oh}}
\def\Spacesstack{\mathcal{S}\!\mathit{paces}}
\def\Quotfunctor{\mathrm{Quot}}
\def\Hilbfunctor{\mathrm{Hilb}}
\def\Curvesstack{\mathcal{C}\!\mathit{urves}}
\def\Polarizedstack{\mathcal{P}\!\mathit{olarized}}
\def\Complexesstack{\mathcal{C}\!\mathit{omplexes}}
% \Pic is the operator that assigns to X its picard group, usage \Pic(X)
% \Picardstack_{X/B} denotes the Picard stack of X over B
% \Picardfunctor_{X/B} denotes the Picard functor of X over B
\def\Pic{\mathop{\mathrm{Pic}}\nolimits}
\def\Picardstack{\mathcal{P}\!\mathit{ic}}
\def\Picardfunctor{\mathrm{Pic}}
\def\Deformationcategory{\mathcal{D}\!\mathit{ef}}
\usepackage{cleveref}
\crefname{originallemma}{original lemma}{original lemmas}
\Crefname{originallemma}{Original lemma}{Original lemmas}
\makeatletter
\newlabel{more-algebra-lemma-map-G-ring-to-completion-regular}{{0AH2}{source}{}{sourcepointer-1}{}}
\newlabel{more-algebra-lemma-map-G-ring-to-completion-regular@cref}{{[originallemma][0][]0AH2}{[1][0][]source}{}{}{}}
\makeatother
\crefname{originalproposition}{original proposition}{original propositions}
\Crefname{originalproposition}{Original proposition}{Original propositions}
\makeatletter
\newlabel{more-algebra-proposition-ubiquity-excellent}{{07QW}{source}{}{sourcepointer-2}{}}
\newlabel{more-algebra-proposition-ubiquity-excellent@cref}{{[originalproposition][0][]07QW}{[1][0][]source}{}{}{}}
\makeatother
\crefname{originalsection}{original section}{original sections}
\Crefname{originalsection}{Original section}{Original sections}
\makeatletter
\newlabel{more-algebra-section-regular}{{07BY}{source}{}{sourcepointer-3}{}}
\newlabel{more-algebra-section-regular@cref}{{[originalsection][0][]07BY}{[1][0][]source}{}{}{}}
\makeatother
\title{The restricted formal polynomial algebra is regular over a rig-closed quotient}
\author{The Stacks Project Authors}
\date{stacks-project:restricted.tex:local-snapshot-20261012; GNU FDL 1.2 or later}
\begin{document}
\maketitle
\noindent\textit{Editorial note.} All delivered mathematical text is the author's own, copied from the pinned Stacks Project chapter file restricted.tex; no statement or proof step was rewritten, reordered or removed.\par\medskip
\noindent\textit{Editorial note.} The counted proof is the author's complete argument: the proof routes regularity through the excellence of the complete local quotient and its polynomial ring, applies the G-ring-to-completion regularity lemma, and composes with the smooth polynomial extension\par\medskip
\noindent\textit{Editorial note.} The label lemma-fibre-regular names the counted statement, and the proof environment immediately following it in the source is delivered verbatim as its proof.\par\medskip
\noindent\textit{Editorial note.} The counted statement is carried by the author's own lemma environment in the chapter source.\par\medskip
\noindent\textit{Allowed standard prerequisites (editorial).}\begin{itemize}
\item complete local Noetherian rings are excellent (More on Algebra, Proposition more-algebra-proposition-ubiquity-excellent) and excellence is preserved by polynomial rings
\item the map from a G-ring to its completion is regular (More on Algebra, Lemma more-algebra-lemma-map-G-ring-to-completion-regular)
\item smooth ring maps are regular and regular ring maps are closed under composition (More on Algebra, Section more-algebra-section-regular)
\end{itemize}\medskip
\ifdefined\hypertarget\hypertarget{sourcepointer-1}{}\fi
\noindent\textit{Source reference.} Original lemma 0AH2: sources/stacks-chow/tags.txt, source lines 13215--13215. Prerequisite pointer invoked at proof line(s) 3630; it names a standard fact of a sibling chapter that the counted argument uses, and every step of the counted proof built on it is retained verbatim in the delivered body.\par\medskip
\ifdefined\hypertarget\hypertarget{sourcepointer-2}{}\fi
\noindent\textit{Source reference.} Original proposition 07QW: sources/stacks-chow/tags.txt, source lines 9846--9846. Prerequisite pointer invoked at proof line(s) 3624; it names a standard fact of a sibling chapter that the counted argument uses, and every step of the counted proof built on it is retained verbatim in the delivered body.\par\medskip
\ifdefined\hypertarget\hypertarget{sourcepointer-3}{}\fi
\noindent\textit{Source reference.} Original section 07BY: sources/stacks-chow/tags.txt, source lines 9358--9358. Prerequisite pointer invoked at proof line(s) 3621; it names a standard fact of a sibling chapter that the counted argument uses, and every step of the counted proof built on it is retained verbatim in the delivered body.\par\medskip
\begin{lemma}
\label{lemma-fibre-regular}
Let $A$ be an adic Noetherian topological ring.
Let $\mathfrak p \subset A$ be a rig-closed prime.
For any $n \geq 1$ the ring map
$$
A/\mathfrak p
\longrightarrow
A\{x_1, \ldots, x_n\} \otimes_A A/\mathfrak p =
A/\mathfrak p\{x_1, \ldots, x_n\}
$$
is regular. In particular, the algebra
$A\{x_1, \ldots, x_n\} \otimes_A \kappa(\mathfrak p)$
is geometrically regular over $\kappa(\mathfrak p)$.
\end{lemma}

\begin{proof}
We will use some fact on regular ring maps the reader can find
in More on Algebra, Section \ref{more-algebra-section-regular}.
Since $A/\mathfrak p$ is a complete local Noetherian ring it
is excellent (More on Algebra, Proposition
\ref{more-algebra-proposition-ubiquity-excellent}).
Hence $A/\mathfrak p[x_1, \ldots, x_n]$ is excellent
(by the same reference). Hence
$A/\mathfrak p[x_1, \ldots, x_n] \to A/\mathfrak p\{x_1, \ldots, x_n\}$
is a regular ring homomorphism by
More on Algebra, Lemma
\ref{more-algebra-lemma-map-G-ring-to-completion-regular}.
Of course $A/\mathfrak p \to A/\mathfrak p[x_1, \ldots, x_n]$
is smooth and hence regular. Since the composition of regular
ring maps is regular the proof is complete.
\end{proof}

\end{document}
